\subsubsection{Trajectory space and invariant measure between sheets}
\label{sec:medida-retorno-hojas}

The toroidal topology of APP fixes the support of cardinal translations
and their winding classes. The correspondence between sheets arises from
another operation on that same support: temporal selection of the
additive or multiplicative evaluation. Its geometric typing is
\[
 \eta_{\rm av}(\Sigma)=\mathscr H_{\rm ar},\qquad
 \eta_{\rm av}(\Pi)=\mathscr H_{\rm vol}.
\]
The first sheet corresponds to the additive reading of the boundary; the second,
to the multiplicative reading of the interior. This map transports mode
labels and preserves the incidence \eqref{eq:fav}. The square of the self-scaling
ratio is determined through the trajectories of that incidence.

\Needspace{8\baselineskip}
\begin{definition}[Memory-one symbolic envelope]
Let
\[
 X_{\rm av}=\{(s_k)_{k\in\mathbb Z}\in
 \{\mathscr H_{\rm ar},\mathscr H_{\rm vol}\}^{\mathbb Z}:
 (F_{\rm av})_{s_k,s_{k+1}}=1\text{ for every }k\}.
\]
Index shift acts on this space. It is the smallest
memory-one symbolic system containing the modal sequences and preserving
their consecutive pairs: it admits the two cross-transitions and
volumetric self-retention. Phase, cardinal orientation, and the trajectory
records remain in the TPK state from which this
envelope is obtained.
\end{definition}

\begin{proposition}[Weighting trajectories between sheets]
In areal--volumetric order, the invariant measure of maximal entropy
of \(X_{\rm av}\) has transition matrix and stationary distribution
\begin{equation}
 P_{\rm av}=\begin{pmatrix}0&1\\\varphi^{-2}&\varphi^{-1}\end{pmatrix},
 \qquad
 (\mu_{\rm ar},\mu_{\rm vol})
 =\frac{(1,\varphi^2)}{1+\varphi^2}.
 \label{eq:medida-hojas-exacta}
\end{equation}
In particular, \(\mu_{\rm ar}/\mu_{\rm vol}=\varphi^{-2}\).
\end{proposition}
\begin{proof}
The already constructed incidence has positive eigenvector
\(r=(1,\varphi)^{\mathsf T}\) and eigenvalue \(\varphi\). The
normalization
\[
 (P_{\rm av})_{ij}
 =\frac{(F_{\rm av})_{ij}r_j}{\varphi r_i}
\]
gives the stated matrix. Its rows sum to one because
\(\varphi^{-2}+\varphi^{-1}=1\). The symmetry of \(F_{\rm av}\)
implies that weights proportional to \(r_i^2\) satisfy detailed
balance; therefore, \(\mu P_{\rm av}=\mu\). Irreducibility and
volumetric self-retention determine a unique stationary distribution.

To verify maximality, every allowed transition satisfies
\(\log(P_{\rm av})_{ij}=\log r_j-\log r_i-\log\varphi\).
Under any invariant measure, the endpoint terms cancel
on average. If \(\nu_i\) are its vertex weights and
\(q_i(j)=\nu(s_1=j\mid s_0=i)\), for \(\nu_i>0\), then
\[
 h_\nu\leq H_\nu(s_1\mid s_0)
 =\log\varphi-\sum_{i:\nu_i>0}\nu_iD(q_i\Vert P_i)
 \leq\log\varphi,
\]
where \(D(q\Vert p)=\sum_jq_j\log(q_j/p_j)\) is the relative
entropy on allowed pairs and \(0\log0=0\) is adopted.
The Markov measure in \eqref{eq:medida-hojas-exacta} attains the bound.
Equality in the first inequality requires the Markov property;
equality in the second requires \(q_i=P_i\). Stationarity and
irreducibility then fix \(\mu\), establishing
uniqueness. This is the Parry measure of the defined envelope.
\end{proof}

The frequency \((1/3,2/3)\) counts the instants of the primitive calendar;
\(\mu\) weights the trajectories of the envelope. Its quadratic ratio
expresses the product of the incoming and outgoing incidence amplitudes.
TPK memory and phase restrictions retain their own
domain; the entropy of the envelope corresponds to the symbolic system
specified in this definition.

\subsubsection{Closure operator and limit of marked returns}
\label{sec:operador-retorno-marcado}

The closure coordinate \(\pi\), generated by the regional reader,
enters after the construction of modal incidence. Let
\(e_{\rm ar}=(1,0)^{\mathsf T}\),
\(e_{\rm vol}=(0,1)^{\mathsf T}\), and
\(P_i=e_ie_i^{\mathsf T}\) be the sheet projectors. The reader's
conditions are preservation of both sheets, unit volumetric
normalization, and an areal mark equal to the closure coordinate. These conditions
determine the positive operator
\begin{equation}
 D_\pi=\pi P_{\rm ar}+P_{\rm vol}
       =\begin{pmatrix}\pi&0\\0&1\end{pmatrix}.
 \label{eq:operador-marca-pi}
\end{equation}
Diagonality expresses preservation of sheets, and the two entries
express their normalizations; this makes explicit the uniqueness
of the operator for the fixed reader.

\Needspace{8\baselineskip}
\begin{proposition}[Ratio of boundary returns]
For \(n\geq1\), the marked evaluation
\begin{equation}
 \Theta_n=
 \frac{\langle e_{\rm ar},D_\pi F_{\rm av}^{\,n}e_{\rm ar}\rangle}
 {\langle e_{\rm vol},F_{\rm av}^{\,n}e_{\rm vol}\rangle}
 =\pi\frac{F_{n-1}}{F_{n+1}}
 \label{eq:retorno-marcado-finito}
\end{equation}
satisfies
\begin{equation}
 \lim_{n\to\infty}\Theta_n
 =\pi\frac{\mu_{\rm ar}}{\mu_{\rm vol}}
 =\frac{\pi}{\varphi^2}.
 \label{eq:retorno-marcado-limite}
\end{equation}
\end{proposition}
\begin{proof}
The diagonal entries of \(F_{\rm av}^{\,n}\) count the paths
of length \(n\) returning to their initial sheet. The already proved power
formula provides \(F_{n-1}\) and \(F_{n+1}\). The operator
\(D_\pi\) multiplies the first count by \(\pi\). Finally,
\(F_{n-1}/F_{n+1}\) is the product of two consecutive ratios
converging to \(\varphi^{-1}\); its limit is
\(\varphi^{-2}\). The identity with the weight ratio follows from
\eqref{eq:medida-hojas-exacta}.
\end{proof}

The closure operator appears only once, in the terminal evaluation
of the return. The transfer \(F_{\rm av}\) remains determined
by the calendar. Information from two readers of the same state is thus
composed: modal incidence determines the quadratic weighting,
and the closure region contributes its coordinate to the boundary.

The composition also preserves refinement. For positive cylinders
\(I=[a,b]\) and \(J=[c,d]\), the image under the reader
\(\mathcal R(x,y)=x/y^2\) is the interval
\[
 \mathcal Q(I,J)=\left[\frac{a}{d^2},\frac{b}{c^2}\right].
\]
Monotonicity in each argument proves that the nested closure
and self-scaling cylinders produce nested intervals for the ratio;
positivity of their limits and contraction of their diameters determine
\(\pi/\varphi^2\). The scalar expression thus retains a chain
of operations compatible with depth, which is applied
next in its two orientations.
